\documentclass[12pt,a4paper]{article}

\usepackage[utf8]{inputenc}
\usepackage[T1]{fontenc}
\usepackage[margin=2.5cm]{geometry}
\usepackage{hyperref}
\usepackage{microtype}
\usepackage{parskip}
\usepackage{lmodern}

\hypersetup{colorlinks=true, urlcolor=blue!60!black}

\pagestyle{empty}

\begin{document}

\begin{flushright}
Kundai Farai Sachikonye \\
Auf der Lichtung 19, 82178 Puchheim, Germany \\
\href{https://github.com/fullscreen-triangle}{github.com/fullscreen-triangle} \\
\texttt{kundai.sachikonye@bitspark.com} \\
(+49) 1626386344 \\[6pt]
\today
\end{flushright}

\vspace{1em}

Prof. Dr. Dierk Niessing \\
Institut für Pharmazeutische Biotechnologie \\
Universität Ulm \\
89081 Ulm

\vspace{1.5em}

\textbf{Re: Akademische:r Beschäftigte:r (m/w/d), Institut für Pharmazeutische
Biotechnologie --- Ref.\ 26067}

\vspace{1em}

Dear Professor Niessing,

I am applying for the position of Akademische:r Beschäftigte:r at your institute
(Ref.\ 26067). The role --- strengthening bioinformatic analysis of
high-throughput datasets and AI tooling, both for research and for the new
Pharmaceutical \& Molecular Biotechnology Master's programme --- aligns closely
with both my training and the work I have been doing full-time for several years.

\textbf{High-throughput omics bioinformatics.}
This is my core. I hold an MSc in Computational Biology (Tübingen), and my work
centres on the analysis of transcriptomics, proteomics, and metabolomics data.
\href{https://github.com/fullscreen-triangle/lavoisier}{Lavoisier} is a
high-throughput mass-spectrometry pipeline for proteomics and metabolomics
(98.9\% NIST accuracy, $>$100\,GB datasets);
\href{https://github.com/fullscreen-triangle/gospel}{Gospel} integrates host
genotype, transcriptomics, and metagenomics, validated against 1000~Genomes,
gnomAD, and UK~Biobank. Detecting pathological features in large biological
datasets, with modern AI tooling, is exactly the work I do.

\textbf{Relevance to the institute's focus.}
Your institute studies gene regulation in neuronal disease using high-throughput
methods to detect pathological abnormalities. I have built disease-state models
in precisely this spirit --- including neurodegeneration (Parkinson's and ataxia,
from wearable and physiological signal data) and tumour disease state --- where
the task is to recover a pathological signature from high-dimensional data. I
would be glad to bring these methods to your projects.

\textbf{Teaching and the new programme.}
I am drawn to the chance to help establish the PMB programme, and I bring direct
teaching experience to it: during my doctoral research at TUM I taught at the
Master's level and supervised student work. Together with the extensive open
documentation and tutorial material I have written for every tool I build, and my
experience explaining high-throughput analysis to interdisciplinary,
non-specialist audiences, this means I can contribute to course design and
delivery from the start. I would welcome the opportunity to design and deliver
bioinformatics coursework, supervise student projects, and bring my own ideas to
a programme being built from the ground up.

\textbf{Background and fit.}
My MSc in Pharmaceutical Biotechnology (HAW Hamburg) maps directly onto the
institute's domain, and I have hands-on wet-lab experience (LC-MS/MS and MALDI
workflows at UKE Hamburg; cell culture and microscopy during my BSc thesis),
though I understand this is advantageous rather than required. English is native
(C2); I have been resident in Germany since 2011 and am available to begin on
1~August~2026.

I would welcome the opportunity to discuss how my omics-bioinformatics and AI
background could support your research and the new Master's programme.

\vspace{1.5em}
Yours sincerely,

\vspace{2.5em}
\textbf{Kundai Farai Sachikonye}

\end{document}
